%%%%%%%%%%%%%%%%%
% This is an sample CV template created using altacv.cls
% (v1.3, 10 May 2020) written by LianTze Lim (liantze@gmail.com). Now compiles with pdfLaTeX, XeLaTeX and LuaLaTeX.
% (v1.6.5b, 27 Jun 2023) forked by Nicolás Omar González Passerino (nicolas.passerino@gmail.com)
%
%% It may be distributed and/or modified under the
%% conditions of the LaTeX Project Public License, either version 1.3
%% of this license or (at your option) any later version.
%% The latest version of this license is in
%%    http://www.latex-project.org/lppl.txt
%% and version 1.3 or later is part of all distributions of LaTeX
%% version 2003/12/01 or later.
%%%%%%%%%%%%%%%%

%% If you need to pass whatever options to xcolor
\PassOptionsToPackage{dvipsnames}{xcolor}

%% If you are using \orcid or academicons
%% icons, make sure you have the academicons
%% option here, and compile with XeLaTeX
%% or LuaLaTeX.
% \documentclass[10pt,a4paper,academicons]{altacv}

%% Use the "normalphoto" option if you want a normal photo instead of cropped to a circle
% \documentclass[10pt,a4paper,normalphoto]{altacv}

%% Fork (before v1.6.5a): CV dark mode toggle enabler to use a inverted color palette.
%% Use the "darkmode" option if you want a color palette used to 
% \documentclass[10pt,a4paper,ragged2e,withhyper,darkmode]{altacv}

\documentclass[10pt,a4paper,ragged2e,withhyper]{altacv}

%% AltaCV uses the fontawesome5 and academicons fonts
%% and packages.
%% See http://texdoc.net/pkg/fontawesome5 and http://texdoc.net/pkg/academicons for full list of symbols. You MUST compile with XeLaTeX or LuaLaTeX if you want to use academicons.

% Change the page layout if you need to
\geometry{left=1.2cm,right=1.2cm,top=1cm,bottom=1cm,columnsep=0.75cm}

% The paracol package lets you typeset columns of text in parallel
\usepackage{paracol}

% Change the font if you want to, depending on whether
% you're using pdflatex or xelatex/lualatex
\ifxetexorluatex
  % If using xelatex or lualatex:
  \setmainfont{Roboto Slab}
  \setsansfont{Lato}
  \renewcommand{\familydefault}{\sfdefault}
\else
  % If using pdflatex:
  \usepackage[rm]{roboto}
  \usepackage[defaultsans]{lato}
  % \usepackage{sourcesanspro}
  \renewcommand{\familydefault}{\sfdefault}
\fi

% Fork (before v1.6.5a): Change the color codes to test your personal variant on any mode
\ifdarkmode%
  \definecolor{PrimaryColor}{HTML}{C69749}
  \definecolor{SecondaryColor}{HTML}{D49B54}
  \definecolor{ThirdColor}{HTML}{1877E8}
  \definecolor{BodyColor}{HTML}{ABABAB}
  \definecolor{EmphasisColor}{HTML}{ABABAB}
  \definecolor{BackgroundColor}{HTML}{191919}
\else%
  \definecolor{PrimaryColor}{HTML}{001F5A}
  \definecolor{SecondaryColor}{HTML}{0039AC}
  \definecolor{ThirdColor}{HTML}{F3890B}
  \definecolor{BodyColor}{HTML}{666666}
  \definecolor{EmphasisColor}{HTML}{2E2E2E}
  \definecolor{BackgroundColor}{HTML}{FFFFFF}
\fi%

\colorlet{name}{PrimaryColor}
\colorlet{tagline}{SecondaryColor}
\colorlet{heading}{PrimaryColor}
\colorlet{headingrule}{ThirdColor}
\colorlet{subheading}{SecondaryColor}
\colorlet{accent}{SecondaryColor}
\colorlet{emphasis}{EmphasisColor}
\colorlet{body}{BodyColor}
\pagecolor{BackgroundColor}

% Change some fonts, if necessary
\renewcommand{\namefont}{\Huge\rmfamily\bfseries}
\renewcommand{\personalinfofont}{\small\bfseries}
\renewcommand{\cvsectionfont}{\LARGE\rmfamily\bfseries}
\renewcommand{\cvsubsectionfont}{\large\bfseries}

% Change the bullets for itemize and rating marker
% for \cvskill if you want to
\renewcommand{\itemmarker}{{\small\textbullet}}
\renewcommand{\ratingmarker}{\faCircle}

%% sample.bib contains your publications
%% \addbibresource{main.bib}

\begin{document}
    \name{Angenor N'Gouandi}
    \tagline{INGÉNIEUR LOGICIEL | QUALITÉ \& DÉVELOPPEMENT}
    %% You can add multiple photos on the left or right
    \photoL{4cm}{angenor}

    \personalinfo{
        \email{angenor99@gmail.com}
        \smallskip
        \phone{+225 05-45-29-28-02}
        \location{Abidjan, Côte d'Ivoire}
        \linkedin{angenor-ngouandi-463708186}
        \homepage{angenor.firebaseapp.com}
        %\github{githubUser}
        %\npm{npmUser}
        %\dev{devtoUser}
        %\medium{nicolasomar}
        %% You MUST add the academicons option to \documentclass, then compile with LuaLaTeX or XeLaTeX, if you want to use \orcid or other academicons commands.
        % \orcid{0000-0000-0000-0000}
        %% You can add your own arbtrary detail with
        %% \printinfo{symbol}{detail}[optional hyperlink prefix]
        % \printinfo{\faPaw}{Hey ho!}[https://example.com/]
        %% Or you can declare your own field with
        %% \NewInfoFiled{fieldname}{symbol}[optional hyperlink prefix] and use it:
        % \NewInfoField{gitlab}{\faGitlab}[https://gitlab.com/]
        % \gitlab{your_id}
    }
    
    \makecvheader
    %% Depending on your tastes, you may want to make fonts of itemize environments slightly smaller
    % \AtBeginEnvironment{itemize}{\small}

    \begin{quote}
    \normalsize Ingénieur logiciel avec plus de 5 ans d'expérience en développement applicatif et assurance qualité. Expérience en mise en place de plateformes d'intégration continue et de tests (Squash-TM, Mantis Bug Tracker, Selenium). Consultant pour l'Organisation Internationale de la Francophonie (institution publique internationale) et formateur technique. Master 2 en Système d'Information et Génie Logiciel.
    \end{quote}
    \vspace{-0.5em}

    %% Set the left/right column width ratio to 6:4.
    \columnratio{0.25}

    % Start a 2-column paracol. Both the left and right columns will automatically
    % break across pages if things get too long.
    \begin{paracol}{2}
        % ----- TECH STACK -----
        \cvsection{COMPÉTENCES}
            \textbf{Qualité / Tests} \\[0.3em]
            \cvtag{Squash-TM}
            \cvtag{Mantis}
            \cvtag{Selenium}
            \vspace{0.5em}

            \textbf{Développement} \\[0.3em]
            \cvtag{Python}
            \cvtag{Laravel}
            \cvtag{VueJs}
            \cvtag{Flask}
            \vspace{0.5em}

            \textbf{Bases de données} \\[0.3em]
            \cvtag{Supabase}
            \cvtag{Firebase}
            \cvtag{MySQL}
            \vspace{0.5em}

            \textbf{Outils} \\[0.3em]
            \cvtag{Git}
            \cvtag{Microsoft Azure}
            \cvtag{Linux}
        % ----- TECH STACK -----

        % ----- POINTS FORTS -----
        \cvsection{Points forts}
            \cvachievement{}{Intégration continue}{Squash-TM, Mantis}
            \cvachievement{}{Institution publique}{OIF, ANSUT}
            \cvachievement{}{Support technique}{5 COP depuis 2021}
        % ----- POINTS FORTS -----

        % ----- LANGUAGES -----
        \cvsection{Langues}
            \cvlang{Français}{Langue maternelle} \\
            \cvlang{Anglais}{Intermédiaire}
        % ----- LANGUAGES -----

        % ----- permis -----
        \cvsection{Permis}
            \cvlang{Catégorie}{B, C, D, E} \\
        % ----- permis -----

        % ----- INTERETS -----
        \cvsection{Intérêts}
            \cvtag{Qualité logicielle}
            \cvtag{Open source}
            \cvtag{Automatisation}
            \cvtag{Pédagogie}
        % ----- INTERETS -----
        
        %% Switch to the right column. This will now automatically move to the second
        %% page if the content is too long.
        \switchcolumn
        
        
        
        % ----- EXPERIENCE -----
        \cvsection{Expériences}

            % \cvevent{Consultant développeur mobile}{Beyima}{20 Juillet 2022 -- Aujourd'hui}{Abidjan, Côte d'Ivoire}
            % \begin{itemize}
            %     \item Définition de l'architecture technique des applications mobiles.
            %     \item Prise de décisions sur les technologies et frameworks à utiliser.
            %     \item Participation active au développement des applications (iOS et/ou Android).
            %     \item Optimisation des performances et des expériences utilisateurs des applications.
            %     \item \textcolor{SecondaryColor}{\textbf{Outils: }} Flutter/Dart, Firebase, Laravel API
            % \end{itemize}
            % \divider
            
            
            
            \cvevent{Expert Informatique}{Organisation Internationale de la Francophonie - Institut de la Francophonie pour le Développement Durable}{Mai 2021 -- Aujourd'hui}{Québec, Canada (Télétravail)}
            \begin{itemize}
                \item Conception, développement et maintenance de plusieurs plateformes web pour une institution publique internationale (5 projets en production).
                \item Vérification de la conformité des applications aux spécifications fonctionnelles et techniques avant chaque mise en production.
                \item Suivi et correction des anomalies en collaboration avec les équipes métier lors des phases de recette.
                \item Intégration d'API tierces (Zoom, YouTube, Gemini, ElevenLabs) avec tests de conformité et de performance.
                \item Développement de l'application mobile du guide du négociateur francophone.
                \item \textcolor{SecondaryColor}{\textbf{Outils: }} Python, Laravel, VueJs, Supabase, Flutter, API REST.
            \end{itemize}
            \divider

            \cvevent{Enseignant vacataire en Informatique}{École Supérieure Africaine des TIC (ESATIC), Ministère de l'enseignement supérieur et de la recherche scientifique}{Octobre 2023 -- Aujourd'hui}{Abidjan, Côte d'Ivoire}
            \begin{itemize}
                \item Formation des étudiants de Licence 3 de la filière Développement d'Applications et Systèmes d'Information (DASI) en {\textbf{Programmation Web et Multimédia}}.
                \item Membre du jury des soutenances de Licence 3.
                \item \textcolor{SecondaryColor}{\textbf{Outils: }} HTML, CSS, PHP, Javascript
            \end{itemize}
            \divider

            \cvevent{Formateur Python}{RMO Capital Humain}{Août 2025}{Abidjan, Côte d'Ivoire}
            \begin{itemize}
                \item Formation de développeurs expérimentés (10+ ans d'expérience) en programmation Python chez RMO, entreprise de ressources humaines comptant plus de 30 000 employés.
            \end{itemize}
            \divider

            

            \cvevent{Formateur Flutter}{Orange Côte d'Ivoire}{Avril 2023}{Abidjan, Côte d'Ivoire}
            \begin{itemize}
                \item Formation de porteurs de projets au développement d'applications mobiles avec Flutter pour Orange Côte d'Ivoire.
                \item \textcolor{SecondaryColor}{\textbf{Outils: }} Dart, Flutter, Firebase
            \end{itemize}
            
            \newpage

        \divider

        \cvevent{Développeur Machine Learning}{Société Ivoirienne d'Intelligence Numérique (SIIN)}{Juin 2020 -- Août 2020}{Abidjan, Côte d'Ivoire}
        \begin{itemize}
            \item Conception et entraînement d'un modèle de classification de texte (NLP) pour la détection automatique de Fake News.
            \item Traitement et préparation des données textuelles (pipeline TF-IDF).
            \item Déploiement de la solution via une API Flask et intégration dans une extension Google Chrome.
            \item \textcolor{SecondaryColor}{\textbf{Outils: }} Python, Scikit-learn, NLP (TF-IDF), Flask, Jupyter Notebook.
        \end{itemize}

        \divider
        
        \cvevent{Développeur / Assurance Qualité}{Agence Nationale Du Service Universel Des Télécommunications (ANSUT)}{Avril 2019 -- Août 2019}{Abidjan, Côte d'Ivoire}
        \begin{itemize}
            \item Mise en place d'une plateforme d'intégration continue et de gestion des tests pour les projets sectoriels de l'ANSUT.
            \item Déploiement et configuration de Squash-TM (gestion des plans et cas de tests) et Mantis Bug Tracker (suivi des anomalies).
            \item Automatisation de tests et de tâches via Python et Selenium IDE.
            \item Participation aux réunions de pilotage de projets avec les chefs de projet.
            \item \textcolor{SecondaryColor}{\textbf{Outils: }} Squash-TM, Mantis Bug Tracker, Selenium IDE, Python, Microsoft Azure, Tomcat, Linux (CentOS).
        \end{itemize}
        % ----- EXPERIENCE -----



        % ----- PROJETS -----
        \cvsection{Projets notables}
            \cvevent{Université Senghor}{\cvreference{137.74.117.231}{http://137.74.117.231/}}{Freelance}{Alexandrie, Égypte}
            \begin{itemize}
                \item Refonte complète de la plateforme web de l'Université Senghor de la Francophonie. Projet récemment achevé — le domaine usenghor-francophonie.org sera connecté prochainement (il pointe actuellement vers l'ancienne plateforme).
            \end{itemize}
            \divider

            \cvevent{Epavillon Climatique}{\cvreference{epavillonclimatique.francophonie.org}{https://epavillonclimatique.francophonie.org/}}{OIF}{Québec, Canada}
            \begin{itemize}
                \item Plateforme de gestion des activités soumises pour le Pavillon de la Francophonie lors des COP.
            \end{itemize}
            \divider

            \cvevent{Le Carré des Études}{\cvreference{lecarredesetudes.com}{https://lecarredesetudes.com/}}{Freelance}{}
            \begin{itemize}
                \item Plateforme numérique de publication de magazines téléchargeables avec tableau de bord d'administration et suivi statistique.
            \end{itemize}
            \divider

            \cvevent{Observatoire des Mines de Madagascar}{\cvreference{miningobs.mg}{https://miningobs.mg/}}{Freelance}{Madagascar}
            \begin{itemize}
                \item Développement du site web de l'Observatoire des Mines de Madagascar.
            \end{itemize}
            \divider

            \newpage

            \cvevent{Suivi des revenus miniers}{\cvreference{kaominina-mangarahara.mg}{https://kaominina-mangarahara.mg/}}{Freelance}{Madagascar}
            \begin{itemize}
                \item Plateforme numérique de suivi de l'utilisation des revenus miniers des collectivités territoriales.
            \end{itemize}
        % ----- PROJETS -----

        % ----- MISSIONS -----
        \cvsection{MISSIONS}
            \cvevent{Conférences des Parties (COP 26 à 30)}{Organisation Internationale de la Francophonie}{2021 -- 2025}{}
            \begin{itemize}
                \item Support technique du Pavillon de la Francophonie lors des COP : Glasgow, Royaume-Uni (2021) — Charm el-Cheikh, Égypte (2022) — Dubaï, Émirats Arabes Unis (2023) — Bakou, Azerbaïdjan (2024) — Belém, Brésil (2025).
            \end{itemize}
            \divider

            \cvevent{Salon des sciences, technologies et innovations environnementales}{Organisation Internationale de la Francophonie}{2024 -- 2025}{}
            \begin{itemize}
                \item Yaoundé, Cameroun (Février 2025) — Kinshasa, RDC (Novembre 2024).
            \end{itemize}
        % ----- MISSIONS -----
        
        % ----- EDUCATION -----
        \cvsection{FORMATIONS}
            \cvevent{Master 2 Système d'Information et Génie Logiciel}{École Supérieure Africaine des TIC (ESATIC)}{2020 -- 2021}{Abidjan, Côte d'Ivoire}
            \divider
            
            \cvevent{Licence en Système Réseaux Informatiques et Télécommunication}{École Supérieure Africaine des TIC (ESATIC)}{2018 -- 2019}{Abidjan, Côte d'Ivoire}
            \divider

            \cvevent{Baccalauréat série C}{Lycée Classique d'Abidjan}{2016}{Abidjan, Côte d'Ivoire}

        % ----- EDUCATION -----
        
    \end{paracol}
\end{document}